% concatenated from Homeo1_MA_update.tex
%Version 3 October 2023
% See section 11 of the User Manual for version history
%
%%%%%%%%%%%%%%%%%%%%%%%%%%%%%%%%%%%%%%%%%%%%%%%%%%%%%%%%%%%%%%%%%%%%%%
%%                                                                 %%
%% Please do not use \input{...} to include other tex files.       %%
%% Submit your LaTeX manuscript as one .tex document.              %%
%%                                                                 %%
%% All additional figures and files should be attached             %%
%% separately and not embedded in the \TeX\ document itself.       %%
%%                                                                 %%
%%%%%%%%%%%%%%%%%%%%%%%%%%%%%%%%%%%%%%%%%%%%%%%%%%%%%%%%%%%%%%%%%%%%%

%%\documentclass[referee,sn-basic]{sn-jnl}% referee option is meant for double line spacing

%%=======================================================%%
%% to print line numbers in the margin use lineno option %%
%%=======================================================%%

%%\documentclass[lineno,sn-basic]{sn-jnl}% Basic Springer Nature Reference Style/Chemistry Reference Style

%%======================================================%%
%% to compile with pdflatex/xelatex use pdflatex option %%
%%======================================================%%

%%\documentclass[pdflatex,sn-basic]{sn-jnl}% Basic Springer Nature Reference Style/Chemistry Reference Style



\documentclass[sn-mathphys-num]{sn-jnl}% Math and Physical Sciences Numbered Reference Style 


%%%% Standard Packages
%%<additional latex packages if required can be included here>

\usepackage{graphicx}%
\usepackage{multirow}%
\usepackage{amsmath,amssymb,amsfonts}%
\usepackage{amsthm}%
\usepackage{mathrsfs}%
\usepackage{multicol}
\usepackage{verbatim}
\usepackage{url}
\usepackage{enumitem}
\usepackage{mathtools}
\usepackage[title]{appendix}%
\usepackage[dvipsnames]{xcolor}%
\usepackage{textcomp}%
\usepackage{manyfoot}%
\usepackage{booktabs}%
\usepackage{algorithm}%
\usepackage{algorithmicx}%
\usepackage{algpseudocode}%
\usepackage{listings}%
%%%%

%%%%%=============================================================================%%%%
%%%%  Remarks: This template is provided to aid authors with the preparation
%%%%  of original research articles intended for submission to journals published 
%%%%  by Springer Nature. The guidance has been prepared in partnership with 
%%%%  production teams to conform to Springer Nature technical requirements. 
%%%%  Editorial and presentation requirements differ among journal portfolios and 
%%%%  research disciplines. You may find sections in this template are irrelevant 
%%%%  to your work and are empowered to omit any such section if allowed by the 
%%%%  journal you intend to submit to. The submission guidelines and policies 
%%%%  of the journal take precedence. A detailed User Manual is available in the 
%%%%  template package for technical guidance.
%%%%%=============================================================================%%%%

%% as per the requirement new theorem styles can be included as shown below
\theoremstyle{thmstyleone}%
\newtheorem{theorem}{Theorem}%  meant for continuous numbers
%%\newtheorem{theorem}{Theorem}[section]% meant for sectionwise numbers
%% optional argument [theorem] produces theorem numbering sequence instead of independent numbers for Proposition
\newtheorem{prop}[theorem]{Proposition}% 
%%\newtheorem{proposition}{Proposition}% to get separate numbers for theorem and proposition etc.
\newtheorem{lemma}{Lemma}[section]
\newtheorem{cor}[lemma]{Corollary}
\newtheorem{obs}[lemma]{Observation}

\theoremstyle{thmstyletwo}%
\newtheorem{example}{Example}%
\newtheorem{remark}{Remark}%

\theoremstyle{thmstylethree}%
\newtheorem{defn}{Definition}%

\DeclareMathOperator{\aut}{Aut}
\DeclareMathOperator{\baut}{\textbf{Aut}}
\DeclareMathOperator{\homeo}{Homeo}
\DeclareMathOperator{\bhomeo}{\textbf{Homeo}}
\DeclareMathOperator{\diff}{Diff}
\DeclareMathOperator{\thomeo}{Homeo^{1}}
\DeclareMathOperator{\bthomeo}{\textbf{Homeo}^{\textbf{1}}}
\DeclareMathOperator{\mcg}{MCG}
\DeclareMathOperator{\supp}{supp}
\DeclareMathOperator{\id}{id}

\raggedbottom
%%\unnumbered% uncomment this for unnumbered level heads

\begin{document}

\title[Article Title]{Homeomorphisms of surfaces that preserve continuously differentiable curves}

%%=============================================================%%
%% GivenName	-> \fnm{Joergen W.}
%% Particle	-> \spfx{van der} -> surname prefix
%% FamilyName	-> \sur{Ploeg}
%% Suffix	-> \sfx{IV}
%% \author*[1,2]{\fnm{Joergen W.} \spfx{van der} \sur{Ploeg} 
%%  \sfx{IV}}\email{iauthor@gmail.com}
%%=============================================================%%

\author[1,2]{\fnm{Katherine Williams} \sur{Booth}}\email{k.wbooth3@gmail.com}


\affil[1]{\orgdiv{School of Mathematics}, \orgname{Georgia Institute of Technology}, \orgaddress{ \city{Atlanta}, \state{GA}, \country{USA}}}

\affil[2]{\orgdiv{Department of Mathematics}, \orgname{Vanderbilt University}, \orgaddress{\city{Nashville}, \state{TN}, \country{USA}}}



\abstract{In this paper, we study $\homeo^1(S)$, the group of homeomorphisms of a surface that preserve the set of one-dimensional $C^1$ submanifolds of that surface. The group $\homeo^1(S)$ belongs to a family of similarly defined groups $\homeo^k(S)$ that were recently introduced by the author. In a separate paper, we have shown that for most closed surfaces, $\homeo^k(S)$ is naturally isomorphic to the automorphisms of a smooth fine curve graph. By contrast, the work in this paper gives local conditions that characterize $\homeo^1(S)$. We show that there exists a collection of conditions that are both necessary and sufficient for a homeomorphism of the surface to be an element of this group. These conditions primarily depend upon the structure of the induced map on the projective tangent bundle. Additionally, we provide examples of several types of elements of $\homeo^1(S)$ that are not diffeomorphisms. These include inducing discontinuous maps on the projective tangent bundle and having infinitely many non-differentiable points. }


%%\keywords{keyword1, Keyword2, Keyword3, Keyword4}

%%\pacs[JEL Classification]{D8, H51}

\pacs[MSC Classification]{57K20, 20F65, 57S05}

\maketitle

\section{Introduction}\label{intro}

Given a smooth surface $S$,  we say that a simple closed curve $\alpha$ in $S$ is $C^1$ if it is a properly embedded one-dimensional $C^1$ submanifold of $S$. Equivalently,  these can also be defined as the image of a continuously differentiable and injective $\gamma: S^1 \rightarrow S$ with $\gamma'(t) \neq 0$ for all $t \in S^1$.  We define the group $\thomeo(S)$ to be the set of $f\in\homeo(S)$ such that both  $f$ and $f^{-1}$ map every $C^1$ curve to a $C^1$ curve.   In a separate paper \cite{BoothAutck}, we show that when $S$ is a compact orientable surface with genus at least 2,  $\thomeo(S)$ is naturally isomorphic to the automorphisms of the $C^1$-curve graph. 

The group of $C^1$ diffeomorphisms of $S$,  denoted $\diff^1(S)$,  is a subgroup of $\thomeo(S)$.  We show in Section~\ref{examples} that this containment is proper by adapting an example of Le Roux--Wolff~\cite{LRW}.  
Since there are homeomorphisms that do not map every $C^1$ curve to a $C^1$ curve, we have the following inclusions:
$$\diff^1(S)\; \subsetneq \;  \thomeo(S) \;  \subsetneq \; \homeo(S)$$ 

In this paper, 
we give several new examples of types of elements from $\thomeo(S) {\setminus} \diff^1(S)$. These examples exhibit increasingly interesting and unexpected behavior.  Our Main Theorem provides sufficient and necessary local conditions for elements of $\homeo(S)$ to be in $\thomeo(S)$.  


The remainder of this introduction is organized as follows.  In Section~\ref{intromainthm}, we define relevant structures and state our Main Theorem.  We also give brief descriptions of our examples,  which are defined in more detail in Section~\ref{examples}.   In Section~\ref{introtangentbundle},  we go into a further discussion on the projective tangent bundle of the surface and its connections to $C^1$ curves.  Section~\ref{introquestions} contains a collection of open questions related to $\thomeo(S)$.  Finally, in Section~\ref{introproofidea},  we give an overview of the proof of our theorem and an outline of the paper.  

\subsection{Statement of our Main Theorem}\label{intromainthm}

Before stating our result, we give a couple of necessary definitions.

\medskip

\noindent \textit{Maps on projective tangent spaces.}
In Section~\ref{sectionforward}, we show that any homeomorphism that maps every $C^1$ curve to a $C^1$ curve preserves the equivalence classes on curves through a point associated to their tangent line at that point.  Such a tangent-preserving homeomorphism $f$ induces well-defined maps $$\bar{d}f_p: \mathbb{P}T_pS \rightarrow \mathbb{P}T_{f(p)}S$$ of the projective tangent spaces for every $p \in S$ and also on the projective tangent bundle $$\bar{d}f:~\mathbb{P}TS \rightarrow \mathbb{P}TS.$$

\noindent We denote points in $\mathbb{P}TS$ by $(p, \ell)$, where $p \in S$ and $\ell \in \mathbb{P}T_pS$. 

\medskip

\noindent \textit{Converging along lines and transverse sequences. } Let $\ell$ be a line with slope $m$ in $\mathbb{R}^2$ through the point $(x, y)$.  A sequence of points $\{(x_n, y_n)\}$ in $\mathbb{R}^2$ is said to \emph{converge along  $\ell$ to $(x, y)$} if 
$$\lim_{n \rightarrow \infty} \frac{y_n - y}{x_n - x} = m$$

To expand this definition to any surface, we identify each tangent space $T_pS$ with $\mathbb{R}^2$.   For any line $\ell \in T_pS$ through $p$,  we say that the sequence $\{p_n\}$ in $S$ \emph{converges along $\ell$ to $p$} if there exists a smooth coordinate chart $\varphi: U \subset S \rightarrow T_pS$ based at $p$ such that $\varphi(p_n)$ converges along $\ell$ to $\varphi(p)$.  Such a chart can always be obtained by fixing a Riemannian metric and taking the inverse of an exponential map based at $p$. 

Any line $\ell \in T_pS$ through the origin is also naturally identified with a point in $\mathbb{P}T_pS$.  We abuse notation and denote this point also by $\ell \in \mathbb{P}T_pS$.  A \emph{transverse sequence for $(p, \ell)$} is a sequence $\{(p_n, k_n)\}$ in $\mathbb{P}TS$ such that $\{p_n\}$  converges along $\ell$ to $p$ and no subsequence of $(p_n, k_n)$ converges to $(p, \ell)$ in $\mathbb{P}TS$.  If such a $(p, \ell)$ exists for a sequence $\{(p_n, k_n)\}$, then we call $\{(p_n, k_n)\}$ a \emph{transverse sequence}.

\medskip

\noindent We now give the necessary and sufficient properties that a homeomorphism needs to have in order to be an element of $\thomeo(S)$.

\begin{theorem}\label{theoremconditions}
Fix a smooth surface $S$.  A homeomorphism $f$ of $S$ is an element of $\thomeo(S)$ if and only if $f$ has the following three properties:
\begin{enumerate}[noitemsep,topsep=0pt, label=$(\alph*)$]
\item $f$ maps every $C^1$ curve to a $C^1$ curve,
\item $\bar{d}f_p$ is a homeomorphism for all $p \in S$, and 
\item $f$ maps every transverse sequence to a transverse sequence.
\end{enumerate}
\end{theorem}

\noindent While $\diff^1(S)$ is a subgroup of $\thomeo(S)$ directly from the definition, it is straightforward to see that any elements of $\diff^1(S)$ also has the three properties stated in the Main Theorem. 
Property $(a)$ follows straight from definitions.
For property~$(b)$, the differential maps between the tangent spaces are linear, so they induce homeomorphisms of projective tangent spaces.  Finally,  property $(c)$ is given by elements of $\diff^1(S)$ acting continuously on the tangent bundle.  

\medskip

\noindent \textit{Example elements.} In addition to the our Main Theorem, we also give examples in Section \ref{examples} of several phenomena that appear in elements of $\thomeo(S)$.   These correspond to the following two results:

\begin{prop}\label{propexamples}
There exist elements $f \in \thomeo(S){\setminus} \diff^1(S)$ and sequences of points $p_n \in S$ such that 
\begin{enumerate}[noitemsep,topsep=0pt, label=$(\roman*)$]
\item $\bar{d}f_{p_1}$ is the identity, but $f$ is not differentiable at $p_1$, 
\item the pushforward map $df_{p_2} : T_{p_2}S \rightarrow T_{f(p_2)}S$ is well-defined and maps every element in $T_{p_2}S$ to~$0$, 
\item $\bar{d}f$ is not continuous, and 
\item $f$ is not differentiable at $p_n$, for all $n \in \mathbb{N}$.
\end{enumerate}
\end{prop}

\begin{prop}\label{propNex}
There exists an element of $\homeo(S){\setminus} \thomeo(S)$ that has properties $(a)$ and $(b)$ from the Main Theorem.
\end{prop}

\noindent In particular, this final result highlights the necessity of property $(c)$ from the Main Theorem.


\subsection{The projective tangent bundle and transverse sequences}\label{introtangentbundle}

In this section, we discuss the projective tangent bundle on a surface and a different way to think about transverse sequences.  Since we are exclusively working with surfaces,  the projective tangent bundle is a circle bundle over the surface.  

\medskip

\noindent \textit{Lifts of $C^1$ curves in the projective tangent bundle.   }  One distinct characteristic of $C^1$ curves is that they have tangent lines at every point and this tangent line varies continuously.  As a consequence,  every $C^1$ curve uniquely defines a continuous curve in the projective tangent bundle.  

But not all continuous curves in the projective tangent bundle are lifts of $C^1$ curves.  One example is a curve in $\mathbb{P}T\mathbb{R}^2$ with points in a neighborhood of the origin in $\mathbb{R}^2$ along the $x$-axis,  but with lines in the vertical direction of the projective tangent spaces for every point.  There is a disconnect between the specified tangent directions and the actual horizontal direction along which the points are approaching the origin.



\begin{figure}[h]
\centering
\includegraphics[width=119mm]{sinewithtransverse.eps}
\caption{A transverse sequence of tangent lines from $y=x^2 \sin(1/x)$ with slope 1 that converge to the origin along the $x$-axis}
\end{figure}

\noindent \textit{Transverse sequences as obstructions to $C^1$ curves.  } In this context, a transverse sequence is an obstruction to the existence of a curve in the projective tangent bundle.  The points $\{(p_n, k_n)\}$ in $\mathbb{P}TS$ corresponding to a transverse sequence cannot converge to the point $(p,\ell)$.  But any differentiable curve in $S$ that passes through infinitely many of the $p_n$ must have a tangent line of $\ell$ at $p$.  So the lift of a curve in $S$ that contains a subsequence of $\{(p_n, k_n)\}$ is not continuous and cannot be a curve in $\mathbb{P}TS$. 

Similarly,  any differentiable curve that is not $C^1$ has tangent lines that do not vary continuously and must contain a transverse sequence.  Thus property $(c)$ from the Main Theorem prevents a such a curve from being mapped to a $C^1$ curve.  This prevents the inverse from sending any $C^1$ curve to a non-$C^1$ curve. 





\subsection{Further questions}\label{introquestions}
The properties of $\thomeo(S)$ and the examples in Section \ref{examples} suggest several possible questions about $\thomeo(S)$.  We have listed a few particular questions below.

\medskip 
\noindent \textbf{Question 1.} Can the set of non-differentiable points in an element of $\thomeo(S)$ be uncountable? Or have positive measure?


\medskip
\noindent \textbf{Question 2.} Is the group $\thomeo(S)$ perfect?

\medskip
\noindent \textbf{Question 3.} Does $\thomeo(S)$ satisfy automatic continuity?

\medskip 

\noindent The previous two questions have been studied for both $\homeo(S)$ and $\diff(S)$. For compact $S$, they have both been answered positively,  but the tools used for $\homeo(S)$ are different from those used for $\diff(S)$.    
Further discussion about these properties and others can be found in an article by Mann~\cite{Mann}. 

\medskip

\noindent \textbf{Question 4.} What is a natural topology on $\homeo^1(S)$? 

\medskip

\noindent  This question was suggested by Belegradek in a personal comment.  Since $\homeo^1(S)$ is a subgroup of $\homeo(S)$, then $\thomeo(S)$ inherits the compact-open topology from $\homeo(S)$.  But this might not be the one most natural one for $\homeo^1(S)$ since it does not account for the additional structure on the projective tangent bundle.  For example, the compact-open topologies for $\diff^1(S)$ also require that derivatives also converge uniformly.   


\subsection{Paper structure and overview of the proof}\label{introproofidea}
Our paper is split into three parts.  The first part is Section~\ref{examples}, where we give several types of examples showing Propositions \ref{propexamples} and \ref{propNex} stated above.  We then prove our Main Theorem.  The proof is split into two parts, each corresponding to one of the logical directions of the statement.  The forward direction involves recovering each of the three properties stated in the theorem from an element of $\thomeo(S)$,  while the reverse direction requires us to show that any map with these properties is an element of $\thomeo(S)$.  

\medskip

\noindent \textit{Forward direction.} We complete this step in Section~\ref{sectionforward}.  
To prove the forward direction, we must have that any element of $\homeo^1(S)$ maps every $C^1$ curve to a $C^1$ curve,  induces well-defined homeomorphisms on the projective tangent space,  and sends transverse sequences to transverse sequences.  

The first property follows directly from the definition of $\thomeo(S)$.  The next step is to show that an element of $\thomeo(S)$ induces a well-defined map on $\mathbb{P}TS$.   We prove this by showing that any homeomorphism of $S$ that sends every $C^1$ curve to a $C^1$ curve must respect the equivalence classes at each point of curves with the same tangent line.

To obtain the second property, we utilize the identification of the projective tangent space with $S^1$ to obtain that the induced map is a homeomorphism. This follows from the fact that bijections of $S^1$ preserving cyclically ordered triples are homeomorphisms.

Finally,  to show that an element of $\thomeo(S)$ sends all transverse sequences to transverse sequences, we work from a contradiction.  We assume that the homeomorphism maps some transverse sequence to a sequence that is not transverse and has a convergent subsequence.  We are then able to construct a $C^1$ curve whose preimage is not $C^1$.

\medskip

\noindent \textit{Reverse direction.}
This step of our proof is completed in Section~\ref{sectionreverse}.  Our goal is to show that any homeomorphism satisfying the three properties from our Main Theorem must be an element of $\thomeo(S)$.  From the first property, we know that $f$ maps every $C^1$ curve to a $C^1$ curve. Thus it only remains to show that $f^{-1}$ also maps every $C^1$ curve to a $C^1$ curve.  

We use the homeomorphism on the tangent space at each point to recover that every image of a $C^1$ curve under $f^{-1}$ must have a unique tangent line at every point.  We then use transverse sequences to show that these tangent lines must vary continuously, which is an equivalent characterization of $C^1$ curves. 


\subsection{Acknowledgments}

We would like to thank our advisor Dan Margalit for his generous support and encouragement.  In particular, we would like to thank him for his frequent suggestions of ways to simplify the description of $\homeo^1(S)$ that gave us the push to find increasingly interesting examples.

We thank Alex Nolte for his relentless enthusiasm and interest in this project.  We especially thank him for his contribution to the arguments involving transverse sequences.

We thank Igor Belegradek for his interest and questions.  We also thank Noah Caplinger, Sean Eli,  Daniel Minahan,  and Roberta Shapiro for helpful conversations and general encouragement for this project. The author was partially supported by the National Science Foundation under Grants No.  DMS-1745583 and DMS-2417920.



\section{Example elements from $\mathbf{\bhomeo^1(S)}$} \label{examples}
In this section, we prove Propositions \ref{propexamples} and \ref{propNex} from Section~\ref{intro}.   We accomplish this by giving several types of maps that are elements of $\thomeo \!\left(\mathbb{R}^2\right){\setminus}\diff^1\!\left(\mathbb{R}^2\right)$.  The primary goal for many of these examples is to dismiss potentially simpler characterizations for elements of $\thomeo(S)$.  Since the properties of $\thomeo(S)$ from our Main Theorem are dependent primarily on the local conditions around a point,  we give these examples explicitly on $\mathbb{R}^2$.  

\subsection{Descriptions of the examples} \label{subexamples}

In this section,  we start with the examples for parts $(i)$ and $(ii)$ from Proposition~\ref{propexamples}.  We then give a non-example that proves Proposition~\ref{propNex}. Finally, we conclude this section with the examples for parts $(iii)$ and $(iv)$ from Propostion~\ref{propexamples}.


\subsection*{Example $\mathbf{(i)}$: Modifying magnitudes with Le~Roux--Wolff} 

This example proves the following proposition related to part $(i)$ from Proposition~\ref{propexamples}.

\begin{prop}\label{propex1}
There exists an element $G \in \thomeo \!\left(\mathbb{R}^2\right){\setminus}\diff^1\!\left(\mathbb{R}^2\right)$ such that $\bar{d}G_{(0,0)}$ is the identity, but $G$ is not differentiable at $(0,0)$. 
\end{prop}

\begin{proof}
We build $G:\mathbb{R}^2 \rightarrow \mathbb{R}^2$ similar to an example given by Le Roux--Wolff \cite[Section 5]{LRW}.  In their example, they change $y$-coordinates based on the slope of the line to the origin.   We modify their construction by modifying both coordinates instead of only one of them.  

Let $g: \mathbb{R} \rightarrow \mathbb{R}^+$ be a smooth map such that $g(1) \neq 1$ and $g(t)=1$ whenever $t$ is outside of the interval $[1/2,2]$.  Define the map $G:\mathbb{R}^2\rightarrow\mathbb{R}^2$ by $$G(x,y) = \left\{ \begin{array}{ll} (x \, g(y/x), y \, g(y/x) ) & \mbox{if } x \neq 0 \\  (x, y) & \mbox{if } x =0 \end{array} \right.$$
Note that $G$ is a $C^1$ diffeomorphism at every point except the origin.  We can show by direct computation that every $C^1$ curve through the origin is mapped by $G$ to a $C^1$ curve.  Since $G^{-1}$ is of the same form as $G$ by replacing $g$ by $1/g$, then $G^{-1}$ also maps every $C^1$ curve to a $C^1$ curve. Thus $G \in \thomeo \!\left(\mathbb{R}^2 \right)$.  

\begin{figure}[h]
\centering
\includegraphics[width=119mm]{LRWmag_together.pdf} 
\caption{The action of $G$ on the vectors $(1, 0)$, $(0, 1)$, and $(1,1)$ based at the origin}
\end{figure}

To show that $G \notin \diff^1 \! \left( \mathbb{R} ^2 \right)$, we consider the action of $G$ on vectors based at the origin.  By changing the magnitude of some of these vectors, we prevent the pushforward map $dG_{(0,0)}$ from being a linear map.   Specifically, $dG_{(0,0)}(1,0) = (1,0)$ and $dG_{(0,0)}(0,1) = (0,1)$,  while $dG_{(0,0)}(1,1) = (g(1), g(1)) \neq (1,1)$.  

Since both $x$ and $y$-coordinates are changed by the same multiplicative factor,  this map does not change the tangent direction of any $C^1$ curve through the origin.  So the induced map $\bar{d}G_{(0,0)}$ is the identity map.  
\end{proof}

\subsection*{Example $\mathbf{(ii)}$: Collapsing and expanding tangent vectors}

For part $(ii)$ from Proposition~\ref{propexamples}, this example proves the following:

\begin{prop}\label{propex2}
There exists an element $H \in \thomeo \!\left(\mathbb{R}^2\right){\setminus}\diff^1\!\left(\mathbb{R}^2\right)$ fixing the origin such that the pushforward map $dH_{(0,0)} : T_{(0,0)}\mathbb{R}^2 \rightarrow T_{(0,0)}\mathbb{R}^2$ is well-defined and maps every element in $T_{(0,0)}\mathbb{R}^2$ to~$0$.
\end{prop}

\begin{proof}
Consider the map $H: \mathbb{R}^2 \rightarrow \mathbb{R}^2$ defined in polar coordinates as follows:
\begin{flalign*}
H: (r,  \theta) & \mapsto (sgn(r) \, r^2, \theta) 
\end{flalign*}

\begin{figure}[h]
\centering
\includegraphics[width=119mm]{Collapse_together.pdf}
\caption{The action of $H$ on a grid near the origin}
\end{figure}


For any unit-speed $C^1$ parameterization of a curve $\gamma(t)$ such that $\gamma(0)=(0,0)$, the curve $H \circ \gamma$ has a tangent vector with magnitude $0$ at the origin.  By reparameterizing by $s=sgn(t)\,  t^2$, it can be shown that the image of $H\circ \gamma$ is still a $C^1$ curve. 

The inverse map
\begin{flalign*}
H^{-1}: (r,  \theta) & \mapsto (sgn(r) \, \sqrt{|r|}, \theta) 
\end{flalign*}

\begin{figure}[h]
\centering
\includegraphics[width=119mm]{Expand_together.pdf}
\caption{The action of $H^{-1}$ on a grid near the origin}
\end{figure}
\noindent acts similarly.  The pushforward map $dH^{-1}_{(0,0)}$ sends every nonzero vector based at the origin to a vector with infinite magnitude.  Again by reparameterizing,  it can be shown that $H^{-1}$ also maps every $C^1$ curve to a $C^1$ curve. Thus $H$ is an element of $\thomeo\! \left(\mathbb{R}^2\right)$.
\end{proof}

The examples of $H$ and $H^{-1}$ both stem from the fact that we are concerned about the image of the curve and not a particular parameterization.  This justifies that the action on $\mathbb{P}TS$ is the appropriate object to consider for elements of $\thomeo(S)$ instead of the pushforward maps.

\subsection*{Proving Proposition~\ref{propNex} : A cautionary non-example for $\mathbf{\bthomeo(S)}$}
We now take a brief respite from examples of $\thomeo\!\left(\mathbb{R}^2\right)$ to give an example of an element from $\homeo\!\left( \mathbb{R}^2 \right){\setminus}\thomeo\!\left(\mathbb{R}^2\right)$.  The map $W$ given below takes every $C^1$ curve to a $C^1$ curve, but the inverse $W^{-1}$ takes some $C^1$ curves to curves that are not $C^1$.  

\begin{proof}[Proof of Proposition~\ref{propNex}]
 To formally define $W$, we start with a $C^1$ diffeomorphism $w: [0,1] \rightarrow [0,1]$ such that $w(0)=0$, $w(1)=1$, $w'(0)=0 = w'(1)$,   and  $w'(t) < \infty$.  For $\theta \in [0,\pi)$, we then define

$$W(r, \theta) = \left\{ \begin{array}{ll}  \left(r, \pi \, w\left( \theta / \pi \right) \, (1- w(|r|)) + \theta \, w(|r|) \right) & \mbox{if } |r|<1 \\
(r, \theta) & \mbox{if } |r| \geq 1
\end{array} \right. $$

\begin{figure}[h]
\centering
\includegraphics[width=119mm]{nonauttogether.pdf}
\caption{The action of $W$ on lines passing through the origin}
\end{figure}

It can be shown that $W$ is a homeomorphism that sends every $C^1$ curve to a $C^1$ curve.  Moreover, it induces well-defined homeomorphisms $\bar{d}W_p$ on the projective tangent spaces at every point $p \in \mathbb{R}^2$.  

But $W^{-1}$ does not map every $C^1$ curve to a $C^1$ curve.  At the origin, $\bar{d}W_{(0,0)}$ compresses the angles near $0$ and $\pi$ in such a way that $W$ sends the curve defined by $y=x^2 \sin (1/x)$ to a $C^1$ curve.  Thus $W$ is not an element of $\thomeo\!\left(\mathbb{R}^2\right)$.
\end{proof}


\subsection*{Example $\mathbf{(iii)}$: Non-continuity on the projective tangent bundle}

We return back to examples of $\thomeo \!\left(\mathbb{R}^2\right)$.  Recall that elements of $\thomeo(S)$ induce a map on $\mathbb{P}TS$.  By property $(b)$ from our Main Theorem, we know that this map on $\mathbb{P}TS$ must be a bijection.  But the following result related to part $(iii)$ of Proposition~\ref{propexamples} shows that this induced map does not have to be continuous.  

\begin{prop}\label{propex3}
There exists an element $Q \in \thomeo \!\left(\mathbb{R}^2\right){\setminus}\diff^1\!\left(\mathbb{R}^2\right)$ such that $\bar{d}Q$ is not continuous.
\end{prop}

\begin{figure}[h!]
\centering
 \raisebox{.35in}{\includegraphics[width=119mm]{nonconttantogether.pdf}} 
\caption{The action of $Q$ on a grid near the origin}
\end{figure}

\begin{proof}
First, take a $C^1$ map $q: \mathbb{R} \rightarrow \mathbb{R}^+$ such that $q(0) = 1, q'(0) = 1$, and $q(t) = 1$ whenever $t \notin [-1,1]$. In addition,  pick such a $q$ so that $q(t) > t \, q'(t)$ for all $t$.  

We then define $Q: \mathbb{R}^2 \rightarrow \mathbb{R}^2$ by
$$Q(x,y) = \left\{ \begin{array}{ll} (x , y \,  q(x/y) ) & \mbox{if } y \neq 0 \\  (x, y) & \mbox{if } y =0 \end{array} \right.$$


As before, we denote elements of $\mathbb{P}T\mathbb{R}^2$ by $(p,  m)$ where $p \in \mathbb{R}^2$ is a point and $m \in \mathbb{P}T_p \mathbb{R}^2$.  For convenience, we let $m$ be the slope when the line is not vertical and $\infty$ when the line is vertical.  

We now consider the sequence $\{((0,1/n), 0)\}_{n \in \mathbb{N}}$ of horizontal lines converging to $((0,0), 0)$ in $\mathbb{P}T_{(0,0)}\mathbb{R}^2$.  After applying the induced map $\bar{d}Q$ to this sequence,  the image will be the sequence $\{((0, 1/n), 1)\}_{n \in \mathbb{N}}$. But $\bar{d}Q$ is the identity on $((0,0), 0)$.  Since the sequence $\{((0, 1/n), 1)\}_{n \in \mathbb{N}}$ does not converge to $((0,0), 0)$, then $\bar{d}Q$ is not continuous.  
\end{proof}


\subsection*{Example $\mathbf{(iv)}$: Countably many non-differentiable points}

A natural question to consider is whether elements of $\thomeo(S)$ can have infinitely many non-differentiable points that accumulate on the surface.  The following result related to part $(iv)$ from Proposition~\ref{propexamples} gives an answer to this question. 

\begin{prop}\label{propex4}
There exists an element $P \in\thomeo \!\left(\mathbb{R}^2\right){\setminus}\diff^1\!\left(\mathbb{R}^2\right)$ such that $P$ is not differentiable at the points $\left(2^{-n}, 0 \right)$, for all $n \in \mathbb{N}$. 
\end{prop}

\begin{proof}
We start with Le Roux--Wolfe's \cite{LRW} example element from $\thomeo \!\left(\mathbb{R}^2\right){\setminus}\diff^1\!\left(\mathbb{R}^2\right)$: 
$$P_0(x,y)= \left\{ \begin{array}{ll} (x,x \, p_0(y/x)) & \mbox{if } x \neq 0 \\  (x, y) & \mbox{if } x =0 \end{array} \right.$$
 where $p_0: \mathbb{R}\rightarrow \mathbb{R}$ is a diffeomorphism such that $p_0(x) = x$ when $x$ is outside $[1/2, 2]$, but $p_0(1) \neq 1$.  Our goal is to make a function that has similar behavior at a sequence of points that converge to the origin, while becoming increasingly nicer as the points make the approach.  To that end, we define a sequence of functions
 $$p_k(x) = \frac{1}{2^k} p_0(x)+\left( 1- \frac{1}{2^k} \right)x$$ Note that this family approaches the identity as $k\rightarrow \infty.$ Let $\varphi_0 :\mathbb{R}^2 \rightarrow [0,1]$ be a smooth bump function such that $\varphi_0(x,y) = 1$ when $||(x,y)||<\frac1 2$ and $\varphi_0(x,y) = 0$ when $||(x,y)|| \geq 1$.  Similarly, define a sequence of functions 
$$\varphi_k(x,y) = \varphi_0\left(2^{k+2}(x-2^{-k}), 2^{k+2}y \right)$$
 We can now define the functions 
$$\rho_n(x, y) = \left(x, \left(x-2^{-n} \right)p_n\left( \frac{ y}{x-2^{-n}} \right) \varphi_n(x,y) + \left(1-\varphi_n(x, y) \right) y \right)$$ 

\begin{figure}[h]
\centering
\includegraphics[width=119mm]{suppP3.pdf}
\caption{The support of $P$}
\end{figure}

Note that each $\rho_n$ has a single non-differentiable point and support contained in the open ball of radius $2^{-(n+2)}$ around the point $x_n = (2^{-n},0)$, which means that these supports are all disjoint.  Define $P: \mathbb{R}^2 \rightarrow \mathbb{R}^2$ by 

$$P(x, y) = \left\{ \begin{array}{ll} \rho_n(x, y) & \mbox{when } (x, y) \in \supp(\rho_n) \\
(x, y) & \mbox{otherwise} \end{array} \right. \\[1em]$$

Any curve that is $C^1$ and passes through only finitely many of the $\supp(\rho_n)$ will continue to be $C^1$ after $P$ is applied since each of the non-differentiable points are isolated.   On the other hand,  it can be shown by direct computation that the image of any $C^1$ curve that passes through infinitely many of the $\supp(\rho_n)$ will still be a $C^1$ curve.  In particular, when the $p_n$ approach the identity,  the distortion on the curve also decreases as it approaches the origin, preserving the continuity of the tangent directions.  Thus $P$ maps every $C^1$ curve to a $C^1$ curve.

Since $P^{-1}$ can be described identically, then $P^{-1}$ also maps every $C^1$ curve to a $C^1$ curve. Thus $P$ is an element of $\thomeo \!\left(\mathbb{R}^2\right)$ that is not differentiable at the points  $(2^{-n},0)$.  
\end{proof}

In addition,  this construction also works when point $x_{\infty} = (0,0)$ is non-differentiable and has support on the region of $\mathbb{R}^2$ where $|y|>|x|$.   So the set of non-differentiable points for elements of $\thomeo(S)$ does not have to be a closed set nor finite.


\subsection{The proof of Proposition \ref{propexamples}}

\begin{proof}
We can add the non-differentiable points in $\mathbb{R}^2$ from the examples to an identity map using smooth local charts. 
More precisely,  let $\varphi_i: U_i \subset S \rightarrow \mathbb{R}^2$ for $i\in\{1, 2, 3, 4\}$ be smooth charts such that the $U_i$ are disjoint.
We define the map $f: S \rightarrow S$ by 
$$f(p) = \left\{ \begin{array}{ll} (\varphi_1^{-1} \circ G \circ \varphi_1)(p) & \mbox{ if } p \in U_1 \\
(\varphi_2^{-1} \circ H \circ \varphi_2)(p) & \mbox{ if } p \in U_2 \\
(\varphi_3^{-1} \circ Q \circ \varphi_3)(p) & \mbox{ if } p \in U_3 \\
(\varphi_4^{-1} \circ P \circ \varphi_4)(p) & \mbox{ if } p \in U_4 \\
p & \mbox{ otherwise}
\end{array} \right.$$

By Propositions \ref{propex1}-\ref{propex4},  $f \in \thomeo(S){\setminus}\diff^1(S)$ and has the desired properties.
\end{proof}

\section{Properties of elements of $\mathbf{\bthomeo(S)}$}\label{sectionforward}

The goal of this section is to prove the following proposition,  which is the forward direction of our Main Theorem.

\begin{prop}\label{propthm2forward}
For a smooth surface $S$, if $f \in \thomeo(S)$,  then $f$ has the following three properties:
\begin{enumerate}[noitemsep,topsep=0pt, label=$(\alph*)$]
\item $f$ maps every $C^1$ curve to a $C^1$ curve,
\item $\bar{d}f_p$ is a homeomorphism for all $p \in S$, and 
\item $f$ maps every transverse sequence to a transverse sequence.
\end{enumerate}
\end{prop}

We accomplish this proof in three steps. First, we show that property $(a)$ implies that $\bar{d}f_p$ is well-defined.  We then obtain property $(b)$,  that this induced map is a homeomorphism. Finally, we prove property $(c)$,  that transverse sequences are preserved.

\subsection{A well-defined map of projective tangent space}
In this first section,  to get that $\bar{d}f_p$ is well-defined, we prove the following result:


\begin{lemma}\label{autfinpreservetangent}
Let $f \in \thomeo(S)$.  Let $\alpha$ and $\beta$ be simple closed $C^1$ curves in $S$ and $p \in \alpha \cap \beta$.  If $\alpha$ and $\beta$ have the same tangent line at $p$ then $f(\alpha)$ and $f(\beta)$ have the same tangent line at $f(p)$.
\end{lemma}

Before we prove the Lemma~\ref{autfinpreservetangent}  we need to discuss intersections of $C^1$ curves and prove relevant lemmas for each type of intersection.

\medskip

\noindent \textit{Characterizing intersections of $C^1$ curves.}  Any intersection between two $C^1$ curves can either be an isolated point or a limit point for a sequence of intersection points.  In the isolated case,  the intersection can be either transverse or tangent.  Further, for a tangent intersection point, the curves can either locally cross each other or stay on the same side that they started.   When the curves stay on the same side, we call this a \emph{one-sided} intersection.  The intersections that are either transverse or cross at the tangency are referred to as \emph{topologically transverse} intersections.

The next few lemmas show that tangencies are preserved by elements of $\thomeo(S)$ in all three of these cases: Lemma~\ref{onesidedtangentline} proves the result for one-sided tangents,  Lemma~\ref{crosstangentline} for topologically transverse tangents, and Corollary~\ref{c1infinitesharetangentline} for limit points of intersection points.

\begin{lemma} \label{onesidedtangentline}
Let $f \in \thomeo(S)$.  Let $\alpha$ and $\beta$ be simple closed $C^1$ curves in $S$ that intersect at an isolated point $p$ with a one-sided intersection. Then $f(\alpha)$ and $f(\beta)$ have the same tangent line at $f(p)$.
\end{lemma}

\begin{proof}
Let $\alpha$ and $\beta$ be $C^1$ curves that intersect at a point $p \in S$ in a one-sided intersection.  We note that $\alpha$ and $\beta$ must be tangent at $p$.  If not, there would exist parameterizations of $\alpha$ and $\beta$ such that the tangent vectors at $p$ are not scalar multiples of each other. But then $\alpha$ and $\beta$ would be transverse at $p$.  Thus the intersection at $p$ would also be topologically transverse and not one-sided. 

Since $f$ is a homeomorphism,  then it will preserve one-sided intersections.  Furthermore,  since $f$ maps $C^1$ curves to $C^1$ curves, then $f(\alpha)$ and $f(\beta)$ are $C^1$ curves with a one-sided intersection at $f(p)$.  Thus $f(\alpha)$ and $f(\beta)$ have the same tangent line at $f(p)$.
\end{proof}

\begin{lemma} \label{crosstangentline}
Let $f \in \thomeo(S)$.  Let $\alpha$ and $\beta$ be simple closed $C^1$ curves in $S$ that intersect tangentially at an isolated point $p$ with a topologically transverse intersection. Then $f(\alpha)$ and $f(\beta)$ have the same tangent line at $f(p)$.
\end{lemma}

\begin{proof}
Let $\alpha$ and $\beta$ be simple closed $C^1$ curves in $S$ with $p \in \alpha \cap \beta$ such that $\alpha$ and $\beta$ have the same tangent line at $p$.  If $\alpha$ and $\beta$ are topologically transverse, then there exist simple closed inessential $C^1$ curves $\gamma$ and $\delta$ such that each pair $\{\alpha, \gamma\}$, $\{\alpha, \delta\}$, $\{\beta, \gamma\}$,  $\{\beta, \delta\}$, $\{\gamma,\delta\}$ intersects at $p$ with a one-sided intersection.  This can be accomplished by following either $\alpha$ or $\beta$ and switching to the other curve at $p$,  and then closing the curve without crossing either $\alpha$ or $\beta$.  
\begin{figure}[h]
\centering
\includegraphics[width=119mm]{sametangentarg.pdf}
\caption{Two topologically transverse $C^1$ curves $\alpha$ and $\beta$ at the point $p$ with the one-sided tangent curves $\gamma$ and $\delta$}
\end{figure}

Thus by Lemma \ref{onesidedtangentline},  $f(\alpha), f(\beta), f(\gamma), $ and $f(\delta)$ must all have the same tangent line at $f(p)$. 
\end{proof}

Now that we have both isolated point cases, we move to the case of a limit point of intersection points.  To accomplish this, we first give a result about $C^1$ curves that intersect infinitely many times in $\mathbb{R}^2$.

\begin{lemma}\label{c1infiniteintersectionssharetangent}
If $\alpha$ and $\beta$ are two simple $C^1$ curves in $\mathbb{R}^2$ that intersect at infinitely many distinct points $x_n$ that converge to the origin, then there are parameterizations of $\alpha(t)$ and $\beta(t)$ of $\alpha$ and $\beta$ such that $\alpha'(0) = \beta'(0) \neq 0$.
\end{lemma}

\begin{proof}
By continuity,  the origin must be a point on $\alpha$.  Let $\alpha(t): \mathbb{R} \rightarrow \mathbb{R}^2$ be a $C^1$ parameterization of $\alpha$  such that $\alpha(0)=(0, 0)$ and $\alpha'(t) \neq 0$.  Up to replacing $t$ by $-t$, we can assume that infinitely many of the $x_n$ occur at positive $t$ values.  Let $x_{n_k}$ be the subsequence of of the $x_n$ that have positive $t$ values on $\alpha(t)$.  Similarly,  there is a $C^1$ parameterization $\beta(s): \mathbb{R} \rightarrow \mathbb{R}^2$ of $\beta$ with $\beta(0)=0$ and $\beta'(s) \neq 0$. Let $t_k$  and $s_k$ be sequences such that $\alpha(t_k) = x_{n_{k}} = \beta(s_k)$.  

Since $\alpha(t)$ and $\beta(s)$ are $C^1$ parameterizations and $t_m \rightarrow 0$ and $s_m \rightarrow 0$, then 
\begin{center}
\begin{minipage}{2.5in}
\begin{flalign*}
\alpha'(0) & = \lim_{t \rightarrow 0} \frac {\alpha(t)}{t} \\
& =  \lim_{k \rightarrow \infty} \frac{\alpha(t_k)}{t_k} \\
& =  \lim_{k \rightarrow \infty} \frac{x_{n_{k}}}{t_k} \\
\end{flalign*} \end{minipage} 
\begin{minipage}{2.5in}
\begin{flalign*}
\beta'(0) & = \lim_{s \rightarrow 0} \frac {\beta(s)}{s} \\
& =  \lim_{k \rightarrow \infty} \frac{\beta(s_k)}{s_k} \\
& =  \lim_{k \rightarrow \infty} \frac{x_{n_{k}}}{s_k} \\
\end{flalign*} \end{minipage}
\end{center}
Thus $\alpha'(0)$ and $\beta'(0)$ must lie on the same line, which is the direction that the $x_{n_{k}}$ approach $0$. Since both $\alpha'(0)$ and $\beta'(0)$ are nonzero, then 
$$\frac{\alpha'(0)}{\beta'(0)} = \lim_{k \rightarrow \infty} \frac{s_k}{t_k} = C \in \mathbb{R} \setminus \{0\}$$ Thus $\beta$ can be reparameterized by $s=Ct$ so that $\beta'(0) = \alpha'(0)$. 
\end{proof}

We can directly apply Lemma~\ref{c1infiniteintersectionssharetangent} to surfaces by using a smooth local chart to get the following:

\begin{cor}\label{c1infinitesharetangentline}
If two simple closed $C^1$ curves $\alpha$ and $\beta$ in $S$ intersect at infinitely many distinct points $p_n$ that converge to $p$, then they have the same tangent line at $p$. 
\end{cor}

Since we have dealt with each case individually, we can now complete this section.

\begin{proof}[Proof of Lemma \ref{autfinpreservetangent}]
Let $\alpha$ and $\beta$ be $C^1$ curves in $S$ with the same tangent line at $p \in \alpha \cap \beta$.  If $p$ is an isolated intersection point with a one-sided intersection, then $f(\alpha)$ and $f(\beta)$ also have a one-sided intersection. So by Lemma \ref{onesidedtangentline},  $f(\alpha)$ and $f(\beta)$ have the same tangent line at $f(p)$.  
If $p$ is an isolated intersection point with a topologically transverse intersection, then $f(\alpha)$ and $f(\beta)$ have the same tangent line at $f(p)$ by Lemma~\ref{crosstangentline}. 

If $p$ is not an isolated intersection point, then there is a sequence of intersection points $\{p_n\} \subset \alpha \cap \beta$ that converge to $p$.  Since $f$ is a homeomorphism, then $f(p_n)$ is a sequence of distinct intersection points of $f(\alpha)$ and $f(\beta)$ that converge to $f(p)$. Thus by Corollary~\ref{c1infinitesharetangentline},  $f(\alpha)$ and $f(\beta)$ must have the same tangent line at $f(p)$. 
\end{proof}

\subsection{Property $\mathbf{(b)}$: Homeomorphism on $\mathbf{\mathbb{P}T_pS}$}
By Lemma~\ref{autfinpreservetangent},  for any $f \in \homeo(S)$ that maps every $C^1$ curve to a $C^1$ curve,  $f$ naturally induces a map
$$\bar{d}f_p: \mathbb{P}T_pS \rightarrow \mathbb{P}T_{f(p)}S$$
between the projective tangent spaces.  The goal of this section is to prove the following lemma related to property $(b)$ from our Main Theorem.

\begin{lemma}\label{projtangentspaceishomeo}
Let $f \in \thomeo(S)$.  Then for all $p \in S$,  the induced map $\bar{d}f_p$ is a homeomorphism.
\end{lemma}


We start with showing that the map is bijective.


\begin{lemma}\label{tangentlinemapwelldefined}
Let $f \in \thomeo(S)$.  Then for all $p \in S$,   the induced map $\bar{d}f_p$ is a bijection.
\end{lemma}

\begin{proof}
Note that for any $f \in \thomeo(S)$, the inverse map $f^{-1}$ must also be in $\thomeo(S)$.  Thus by Lemma~\ref{autfinpreservetangent},  $f^{-1}$ induces a map $$\bar{d}f^{-1}_{f(p)}: \mathbb{P}T_{f(p)}S \rightarrow \mathbb{P}T_pS$$ Moreover, since $f \circ f^{-1} = f^{-1} \circ f = \id$, then 
$$
\bar{d}f_p \circ \bar{d}f^{-1}_{f(p)}  = \bar{d} \id_{f(p)} = \id  $$
and 
  $$\bar{d}f^{-1}_{f(p)} \circ \bar{d}f_p  = \bar{d} \id_p = \id $$
Thus $\bar{d}f_p$ is a bijection. 
\end{proof}
 

\noindent To show that $\bar{d}f_p$ is a homeomorphism, we first
observe the following:
\begin{obs}\label{bijwithorderonRishomeo}
If $h: \mathbb{R} \rightarrow \mathbb{R}$ is a bijection that preserves (or reverses) the natural ordering on $\mathbb{R}$,  then $h$ is a homeomorphism of $\mathbb{R}$. 
\end{obs}


\noindent We now use this observation to prove a general result for bijections of $S^1$. 


\begin{lemma}\label{bijwithorderonS1ishomeo}
If $h: S^1 \rightarrow S^1$ is a bijection that preserves (or reverses) the cyclic order of every triple, then $h$ is a homeomorphism.
\end{lemma}

\begin{proof}
By postcomposing with a rotation, we can assume that $h$ fixes some point $c \in S^1$.  Let $\varphi: S^1 \setminus\{c\} \rightarrow \mathbb{R}$ be a homeomorphism such that for any $x, y \in \mathbb{R}$ such that $x<y$,  the cyclic triple $[\varphi^{-1}(x) \varphi^{-1}(y) c]$ of $S^1$ is in clockwise order.   

\begin{figure}[h]
\centering
\includegraphics[width=119mm]{orderingontangent_S1toRmap.pdf} 
\caption{The homeomorphism $\varphi$ from $S^1 \setminus \{c\}$ to $\mathbb{R}$}
\end{figure}


Then $\varphi \circ h|_{S^1 \setminus\{c\}} \circ \varphi^{-1}$ is an order preserving (or reversing) bijection on $\mathbb{R}$. So by Observation \ref{bijwithorderonRishomeo}, $\varphi \circ h|_{S^1 \setminus\{c\}} \circ \varphi^{-1}$ is a homeomorphism. Thus $h|_{S^1 \setminus\{c\}}$ is also a homeomorphism.  Moreover, since any neighborhood of $c$ will be sent by $h|_{S^1 \setminus\{c\}}$ to a neighborhood of $c$, then $h|_{S^1 \setminus\{c\}}$ can be extended to $h$ and remain a homeomorphism on all of $S^1$. 
\end{proof}

The next lemma shows that the induced map on the projective tangent space satisfies the conditions of the previous lemma.

\begin{lemma}\label{homeo1preservetriple}
If $f \in \homeo^1(S)$ is locally orientation preserving (or reversing) at $p \in S$, then the induced map $\bar{d}f_p$ preserves (or reverses) the cyclic order of every triple.
\end{lemma}

\begin{proof}
It suffices to prove this result for the orientation preserving case, since the reasoning is identical in the orientation reversing case, with the corresponding words replaced.

Let $[abc]$ be a cyclic triple of $\mathbb{P}T_pS$ in clockwise order.  Since $a$, $b$, and $c$ are distinct, then there are representative curves $\alpha$, $\beta$, and $\gamma$ that are transverse at $p$ and a connected neighborhood $N$ of $p$ such that these curves are disjoint in $N \setminus \{p\}$.  Denote these disjoint arcs as $\alpha_1$,  $\alpha_2$, $\beta_1$, $\beta_2$, $\gamma_1$, and $\gamma_2$ such that in clockwise order around $p$, these  arcs occur as the cycle $[\alpha_1 \beta_1 \gamma_1 \alpha_2 \beta_2 \gamma_2 ]$.

\begin{figure}[h]
\centering
\includegraphics[width=119mm]{triple_projtanandsurface.pdf} 
\caption{The correspondence between cyclic triples in $\mathbb{P}T_pS$ and their representatives on $S$ near the point $p$}
\end{figure}


Since $f$ is a homeomorphism, then the arcs $f(\alpha_1), f(\alpha_2), f(\beta_1), f(\beta_2), (\gamma_1)$, and $f(\gamma_2)$ are also disjoint in $f(N \setminus\{p\})$.  Moreover,  the cycle $[f(\alpha_1) f(\beta_1) f(\gamma_1) f(\alpha_2) f(\beta_2) f(\gamma_2) ]$ is in clockwise order since $f$ is orientation preserving.  Thus $f(\alpha)$, $f(\beta)$, and $f(\gamma)$ are still topologically transverse at $p$ with the same ordering. 

By Lemma \ref{tangentlinemapwelldefined},  we know that $\bar{d}f_p$ is a bijection. Thus $\bar{d}f_p(a), \bar{d}f_p(b),$ and $\bar{d}f_p(c)$ are distinct and the cyclic triple $[\bar{d}f_p(a) \,\bar{d}f_p(b) \, \bar{d}f_p(c)]$ is in clockwise order. Thus $\bar{d}f_p$ preserves the cyclic order of every triple.
\end{proof}

We can now combine the above results to prove the main lemma of this section.

\begin{proof}[Proof of Lemma \ref{projtangentspaceishomeo}]
By Lemmas \ref{tangentlinemapwelldefined} and \ref{homeo1preservetriple}, we have that for any $f \in \homeo^1(S)$, the map $\bar{d}f_p$ satisfies the conditions of Lemma \ref{bijwithorderonS1ishomeo}.  Thus $\bar{d}f_p$ is a homeomorphism for every $p \in S$. 
\end{proof}

\subsection{Property $\mathbf{(c)}$: Mapping transverse sequences}

In this section, we prove the following lemma, showing  property $(c)$ from our Main Theorem.

\begin{lemma}\label{preservetransverse}
Let $f \in \thomeo(S)$. Suppose that $\{(p_n, \ell_n)\}$ is a transverse sequence that converges to a point $p$ along $\ell$. Then $\{(f(p_n), \bar{d}f_p(\ell_n))\}$ is a transverse sequence that converges to the point $f(p)$ along $\bar{d}f_p(\ell)$.
\end{lemma}

Before we can prove Lemma~\ref{preservetransverse},  we need to give some intermediary results giving the existence of $C^1$ curves through non-transverse sequences in $S \times \mathbb{P}TS$.




\begin{lemma}\label{c1throughtangents}
Let $\{p_n\}$ be a sequence of points in $\mathbb{R}^2 \setminus \{p\}$ that converge along a line $\ell$ to a point $p$.  If $\{(p_n, \ell_n)\}$ is a sequence in $\mathbb{R}^2 \times \mathbb{P}T_{p_n}\mathbb{R}^2$ that converges to $(p,  \ell)$, then there is a subsequence $\{(p_{n_k}, \ell_{n_k})\}$ and a simple closed $C^1$ curve $\gamma$ such that for all $k$, $p_{n_k}$ is a point on $\gamma$ and the tangent line of $\gamma$ at $p_{n_k}$ is $\ell_{n_k}$.
\end{lemma}


\begin{proof}
We complete this proof in two steps. We first find the subsequence of points $\{p_{n_k}\}$ and then we find the $C^1$ curve $\gamma$ that passes through this subsequence.  This construction is inspired by a similar proof from Rosenthal~\cite[Theorem~2]{Rosenthal} that finds a $C^1$ curve through a bounded sequence of points.  We have made adjustments to their argument to account for the additional restriction on tangent lines. 

\medskip

\noindent \textit{Step 1: Finding the subsequence.  } Our goal in this step is to restrict the sequence of converging points in such a way to make it possible to draw a $C^1$ curve through the subsequence.

By composing with an affine map, we can assume that $p$ is the origin and the line $\ell$ is the $x$-axis. 
Since the $p_n$ converge to the origin along the $x$-axis, then only finitely many of them can be on the $y$-axis. So there is a subsequence of $\{p_n\}$ such that either the $x$-values are all positive or all negative.   Additionally, there will be a further subsequence such that the $y$-values are all nonnegative or all nonpositive.  Without loss of generality, we will assume that the subsequence has positive $x$-values and nonnegative $y$-values.  We can further take a subsequence $\{p_{n_s}\}$ such that $|p_{n_s}| > |p_{n_{s+1}}| >0$ for all $s$.  For simplicity of notation, we set $p_s = p_{n_s}$.

\begin{comment}
\begin{figure}[h]
\centering
\includegraphics[width=4in]{c1ptsandtangent.pdf}
\caption{A restricted subsequence $\{(p_{n_k}, \ell_{n_k})\}$ with a $C^1$ arc drawn through the points $p_{n_k}$ with tangent lines $\ell_{n_k}$}
\end{figure}
\end{comment}

Denote by $(ab)$ the slope of the line between the points $a$ and $b$.  Since $\{p_s\}$ is converging to the origin along the $x$-axis, then $(p_s p)$ is also converging to 0.  If there is a subsequence with $(p_{s_k} p ) = 0$, then we take this subsequence. Otherwise we can take a subsequence $\{p_{s_k}\}$ such that $(p_{s_k} p ) > ( p_{s_{k+1}} p) > 0 $.  Furthermore,  note that $(p_{s_k} p_{s_{k+r}}) > (p_{s_k} p)$ for all $r \in \mathbb{Z}^+$. But as $r \rightarrow \infty$ this slope will approach $(p_{s_k} p)$. So we can take another subsequence such that $(p_{s_k} p_{s_{k+1}}) < (p_{s_{k-1}} p_{s_{k}})$ for all $k$. 

Finally, we can take a further subsequence so that the magnitudes of the slopes of $\ell_{s_k}$ converge monotonically to 0.  We now set this final subsequence to be the $\{(p_{n_k}, \ell_{n_k})\}$ that we are looking for. 

\medskip

\noindent \textit{Step 2: Drawing the $C^1$ curve.  } With the restricted sequence $\{(p_{n_k}, \ell_{n_k})\}$, we can now find the $C^1$ curve $\gamma$ that passes through the points $p_{n_k}$ with the tangent lines $\ell_{n_k}$. This step is done by giving exact formulas for several arcs and directly checking that the resulting curve is $C^1$.  

For simplicity,  set $p_{n_k} = (x_k, y_k)$ and let $m_k$ be the slope of $\ell_{n_k}$.  Set $(x_0, y_0) = (x_1+1, y_1)$ and $m_0 =0$.  Note that by choice of subsequence,  $0< x_{k+1} < x_k$ for all $k \in \mathbb{Z}_{\geq 0}$. Then for $t \in [x_{k+1}, x_k]$,  we define the curve 
\begin{flalign*}
\gamma_k(t) & = \left\{ \begin{array} {ll}
\left( \frac{2t- x_k-x_{k+1}}{x_k - x_{k+1}} \right) \left( \left( \frac{(x_k-x_{k+1})m_k}{2 \pi} \right) \, \sin \left( \frac{2 \pi (t-x_k)}{x_k - x_{k+1} } \right)+ y_k \right) & \\
 \qquad + \left( 1 - \left( \frac{2t- x_k-x_{k+1}}{x_k - x_{k+1}}  \right)\right) \left( \frac{y_k - y_{k+1}}{x_k - x_{k+1}} (t-x_{k}) + y_k \right) & \mbox{ if } \frac{x_k+x_{k+1}}{2} \leq t \leq x_k \\
 \left( \frac{2t- x_k-x_{k+1}}{ x_{k+1}-x_k} \right) \left( \left( \frac{(x_k-x_{k+1})m_{k+1}}{2 \pi} \right) \, \sin \left( \frac{2 \pi (t-x_{k+1})}{x_k - x_{k+1} } \right)+ y_{k+1} \right) & \\
 \qquad + \left( 1 - \left( \frac{2t- x_k-x_{k+1}}{x_{k+1}-x_k}  \right)\right) \left( \frac{y_k - y_{k+1}}{x_k - x_{k+1}} (t-x_{k+1}) + y_{k+1} \right) & \mbox{ if } x_{k+1} \leq t \leq  \frac{x_k+x_{k+1}}{2}
\end{array} \right.
\end{flalign*}
Note that 
$$\gamma_k(x_k) = y_k,  \; \gamma_k(x_{k+1}) = y_{k+1},  \; \mbox{ and } \; \gamma_k\left( \frac{x_k+x_{k+1}}{2} \right) = \frac{y_k+y_{k+1}}{2} $$
Moreover, 
\begin{flalign*}
\gamma_k'(t) = \left\{ \begin{array} {ll}
 m_k \left( \left( \frac{2t - x_k - x_{k+1}}{x_k - x_{k+1}} \right)  \cos\left( \frac{ 2 \pi (t-x_k)}{x_k - x_{k+1}} \right) + \frac 1 {\pi} \sin\left( \frac{ 2 \pi (t-x_k)}{x_k - x_{k+1}} \right) \right) & \\ \qquad + \; 2 \left( 1- \frac{2t - x_k - x_{k+1}}{x_k - x_{k+1}} \right) \left( \frac{ y_k - y_{k+1}}{x_k - x_{k+1}} \right) & \mbox{ if } \frac{x_k+x_{k+1}}{2} \leq t \leq x_k \\
  m_{k+1} \left( \left( \frac{2t - x_k - x_{k+1}}{x_{k+1} - x_{k}} \right)  \cos\left( \frac{ 2 \pi (t-x_{k+1})}{x_k - x_{k+1}} \right) - \frac 1 {\pi} \sin\left( \frac{ 2 \pi (t-x_{k+1})}{x_k - x_{k+1}} \right) \right) & \\ \qquad + \; 2 \left( 1- \frac{2t - x_k - x_{k+1}}{x_{k+1} - x_{k}} \right) \left( \frac{ y_k - y_{k+1}}{x_k - x_{k+1}} \right) & \mbox{ if } x_{k+1} \leq t \leq  \frac{x_k+x_{k+1}}{2} \\
\end{array} \right.
\end{flalign*}
with $$\gamma_k'(x_k) = m_k,  \; \gamma_k'(x_{k+1}) = m_{k+1},  \; \mbox{ and } \; \gamma_k'\left( \frac{x_k+x_{k+1}}{2} \right) = 2\frac{y_k-y_{k+1}}{x_k - x_{k+1}} $$
We can see that each $\gamma_k'$ is continuous on its domain. Since both $m_k$ and $\frac{y_k-y_{k+1}}{x_k - x_{k+1}} $ approach 0 as $k \rightarrow \infty$ and all other terms are bounded, then $\gamma_k'(t) \rightarrow 0$ as $k \rightarrow \infty$.  We then define the arc $\xi : [0,x_0] \rightarrow \mathbb{R}^2$ as 
$$\xi(t) = \left\{ \begin{array} {ll}
(0,0) & \mbox{ if } t = 0 \\
\left( t, \gamma_k\left( t \right) \right)& \mbox{ if } x_{k+1} \leq t \leq x_k
\end{array} \right. $$

\noindent Note that $\gamma_k(x_{k+1}) = y_{k+1} = \gamma_{k+1}(x_{k+1})$ and $\gamma'_k(x_{k+1}) = m_{k+1} = \gamma'_{k+1}(x_{k+1})$ and each $\gamma_k$ is $C^1$.  Since $\xi'(t) = (1, \gamma'_k(t)) \neq 0$, then $\xi$ is a $C^1$ arc. 

Now let $ y_{\max} = 1 + \max_{k, t} \gamma_k(t)$.  We then define arcs $\alpha:[0,1] \rightarrow \mathbb{R}^2$,  $\beta: [0,1] \rightarrow \mathbb{R}^2$, and $\delta: [0,1] \rightarrow \mathbb{R}^2$ by:
\begin{flalign*}
\alpha(t) & = \left( y_{\max}  - \frac{ y_0}{2} \right) \left( \cos( \pi ( t + 1/2)),   1+\sin( \pi(t+1/2)) \right) \\
\beta(t) & = \left(x_0(1-t),   2y_{\max}  -  y_0 \right) \\
\delta(t) & = \left( (y_{max} -  y_0) \cos(\pi (t-1/2)) + x_0,   (y_{max} -  y_0) \sin(\pi (t-1/2)) + y_{\max} \right)
\end{flalign*}
Note that $\xi, \alpha, \beta$, and $\delta$ together form a single closed curve $\gamma$ since
\begin{flalign*}
\xi(0) = \alpha(1) & = (0,0) \\
\alpha(0) = \beta(1) & = \left(0,  2y_{\max}  -  y_0 \right)  \\
\beta(0) = \delta(1) & = \left(x_0,  2y_{\max}  -  y_0 \right)  \\
\delta(0) = \xi(x_0) & = (x_0, y_0) 
\end{flalign*}

\begin{figure}[h]
\centering
\includegraphics[width=119mm]{closedc1curve2.pdf} 
\caption{The closed $C^1$ curve $\gamma$ in $\mathbb{R}^2$}
\end{figure}
\noindent In addition,  $\alpha, \beta$ and $\delta$ are all $C^1$ since 
\begin{flalign*}
\alpha'(t) & = \left( y_{\max}  - \frac{ y_0}{2} \right) \left( - \pi \sin( \pi ( t + 1/2)),   \pi \cos( \pi(t+1/2)) \right)  \\
\beta'(t) & = \left( - x_0, 0 \right) \\
\delta'(t) & =  \left( - \pi  (y_{max} -  y_0) \sin(\pi (t-1/2)),  \pi  (y_{max} -  y_0) \cos(\pi (t-1/2)) \right)
\end{flalign*}

\noindent In order to get that the closed curve $\gamma$ is $C^1$, we also need to check is that the tangent lines agree where the arcs $\xi, \alpha, \beta, $ and $\delta$ meet.  We denote the slope of the tangent lines of an arc $\rho(t) = \left(\rho_x(t), \rho_y(t)\right)$ at $t$ by $$m_{\rho}(t) = \frac{\rho_y'(t)}{\rho_x'(t)}$$
It can be directly computed that $$ m_{\xi}(0) = m_{\xi}(x_0) =  m_{\alpha}(0)  = m_{\alpha}(1) = m_{\beta}(0) = m_{\beta}(1) = m_{\delta}(0) = m_{\delta}(1) = 0$$
Since the slopes are the same,  the tangent lines agree at the endpoints of every pair of arcs.  Finally,  we need to show that $\gamma$ is a simple curve.  By construction,  
\begin{flalign*} 
\xi & \subset [0,x_0] \times (-\infty, y_{\max}-1] \\
\alpha & \subset (-\infty, 0] \times (-\infty, \infty)\\
\beta & \subset [0,x_0] \times [y_{\max}, \infty) \\
\delta & \subset [x_0, \infty] \times (\infty, \infty)
\end{flalign*}
The interiors of the arcs are contained in the interiors of these regions, which are all disjoint. So the arcs can only overlap at the endpoints.  

Thus $\gamma$ is a simple closed $C^1$ curve through the sequence of points $\{p_{n_k}\}$ with tangent line $\ell_{n_k}$ at each $p_{n_k}$.  \end{proof}

\noindent By using a smooth chart, we can extend  Lemma~\ref{c1throughtangents} from $\mathbb{R}^2$ to any surface $S$. This gives us the following result:

\begin{cor}\label{c1curvewithtangentsonS}
Let $S$ be a surface and let $\{p_n\}$ be a sequence of points in $S \setminus \{p\}$ that converge along a line $\ell \in \mathbb{P}T_pS$ to a point $p$.  If $(p_n, \ell_n)$ is a sequence in $S \times \mathbb{P}T_{p_n}\mathbb{R}^2$ that converges to $(p,  \ell)$, then there is a subsequence $\{p_{n_k}\}$ and a simple closed $C^1$ curve $\gamma$ such that for all $k$, $p_{n_k}$ is a point on $\gamma$ and the tangent line of $\gamma$ at $p_{n_k}$ is $\ell_{n_k}$.
\end{cor}

\medskip

\noindent \textit{Tangent lines and their angles.  } In $\mathbb{R}^2$, any element of $\mathbb{P}T_{(x, y)}\mathbb{R}^2$ can be naturally identified with the non-obtuse angle between a straight line representative and the horizontal line through $(x, y)$.  We let $\theta_{\ell} \in S^1$ denote the angle corresponding to the line $\ell \in \mathbb{P}T_{(x, y)}\mathbb{R}^2$, with the natural identification $\theta \sim \theta + \pi n$ for $n \in \mathbb{N}$.

Our next lemma gives a characterization  for points converging along a line that utilizes $C^1$ curves.

\begin{lemma}\label{convergealongline}
Let $\{(x_n, y_n)\}$ be a sequence converging to the point $(x, y) \in \mathbb{R}^2$.  Let $\ell \in \mathbb{P}T_{(x, y)}\mathbb{R}^2$.  Then $\{(x_n, y_n)\}$ converges along $\ell$ to $(x,y)$ if and only if for any $\varepsilon >0$  and any $C^1$ curves $\gamma^+$ and $\gamma^-$ whose tangent lines at $(x, y)$ have angles in $(\theta_{\ell}, \theta_{\ell}+ \varepsilon)$ and $(\theta_{\ell}- \varepsilon, \theta_{\ell})$,   there exists a neighborhood $U$ of $(x, y)$ and $N \in \mathbb{N}$ such that for all $n >N$,  $(x_n, y_n)$ is in the same component in $U \setminus (\gamma^+ \cup \gamma^-)$ as the straight line representative of $\ell$ at $(x,y)$.  
\end{lemma}


\begin{figure}[h]
\centering
\includegraphics[width=119mm]{convergingtolinebetweencurves.pdf} 
\caption{The points $(x_n, y_n)$ converging along the line $\ell$ eventually stay between the regions defined by $(-\varepsilon - \delta/2, - \delta / 2) \cup (\delta / 2, \varepsilon + \delta / 2)$ and thus also the curves $\gamma^+$ and $\gamma^-$}
\end{figure}

\begin{proof}
We first prove the forward direction.  After composing by an isometry of $\mathbb{R}^2$, we can assume that $(x,y)$ is the origin and the $x$-axis is the straight line representative of $\ell$. Let $\varepsilon >0$.  Pick $\gamma^+$ and $\gamma^-$ to be any $C^1$ curves with tangent lines at the origin with angles $\theta^+ \in (0,  \varepsilon)$ and  $\theta^- \in (- \varepsilon, 0)$.  


Set $\delta = \min\{|\theta^+|, |\theta^-|\}$.  Note that $\theta^+ \in [\delta, \varepsilon)$ and $\theta^- \in (-\varepsilon, -\delta]$. Since the angles of the tangent lines of $\gamma^+$ and $\gamma^-$ are continuous,  then there is a neighborhood $U$ of the origin such that the curves $\gamma^+$ and $\gamma^-$ stay within the region defined by the angles $(- \varepsilon -\delta/2, -\delta/2) \cup (\delta/2, \varepsilon + \delta/2)$.  

Let $\theta_n$ denote the angle of the point $(x_n, y_n)$ relative to the origin.  Since the $(x_n, y_n)$ converge along the $x$-axis, then $\theta_n$ converge to 0.  So there exists  $N_1 \in \mathbb{N}$ such that for all $n >N_1$, $\theta_n \in (-\delta/2, \delta/2)$.  Moreover, since $(x_n, y_n)$ converge to the origin, then there exists $N_2 \in \mathbb{N}$ such that for all $n >N_2$, the points $(x_n, y_n)$ are contained in $U$.  By setting $N = \max\{N_1, N_2\}$, we have found the desired $U$ and $N$.

We now prove the reverse direction.  Let $\gamma^+$ and $\gamma^-$ be the straight lines with angles $\varepsilon$ and $-\varepsilon$. Since $\varepsilon$ is arbitrary, then the angle of the points $(x_n, y_n)$ relative to the origin must be converging to 0.  Thus the sequence $\{(x_n, y_n)\}$ is converging along the $x$-axis to the origin.  
\end{proof}

Note that Lemma~\ref{convergealongline} also applies locally to sequences that converge along lines to points in a surface $S$ by replacing straight line representatives with locally geodesic representatives.  We can now prove Lemma~\ref{preservetransverse} and Propostion~\ref{propthm2forward}, giving the forward direction of our Main Theorem. 

\begin{proof}[Proof of Lemma~\ref{preservetransverse}]
Let $f \in \thomeo(S)$.  We proceed by contradiction.  Assume that we have a transverse sequence $\{(p_n, \ell_n)\}$ that converges to $p$ along $\ell$ such that $\{(f(p_n), \bar{d}f_p(\ell_n))\}$ is not a transverse sequence.  

Since $f$ is a homeomorphism, then $\{f(p_n)\}$ must converge to $f(p)$.  Next, let $\varepsilon >0$ and consider a pair of curves $\gamma^+$ and $\gamma^-$ from Lemma~\ref{convergealongline} through the point $f(p)$ such that the angles $\theta^+$ and $\theta^-$ are in the interval $(\theta_{\bar{d}f_p(\ell)} - \varepsilon,  \theta_{\bar{d}f_p(\ell)} + \varepsilon)$.   By Lemma~\ref{projtangentspaceishomeo}, the curves $f^{-1}(\gamma^+)$ and $f^{-1}(\gamma^-)$ have tangent lines at $p$ with angles in the intervals $(\theta_{\ell},  \theta_{\ell} + \eta)$ and $(\theta_{\ell} -\eta, \theta_{\ell})$ respectively, for some $\eta >0$.   Since $\{p_n\}$ converges along $\ell$ to $p$,  there is a neighborhood of $U$ of $p$ and $N \in \mathbb{N}$ such that for all $n >N$, the points $p_n$ are in the same component of $U \setminus (f^{-1}(\gamma^+) \cup f^{-1}(\gamma^-))$ as a locally geodesic representative of $\ell$.  But then for all $n >N$,  the points $f(p_n)$ are in the same component of $f(U) \setminus (\gamma^+ \cup \gamma^-)$ as a locally geodesic representative of $\bar{d}f_p(\ell)$.  Thus by Lemma~\ref{convergealongline},  the sequence $\{f(p_n)\}$ converges along $\bar{d}f_p(\ell)$ to $f(p)$. 

For $\{(f(p_n), \bar{d}f_p(\ell_n))\}$ to not be a transverse sequence,  it must have a subsequence that converges to $(f(p), \bar{d}f_p(\ell))$.  By Corollary~\ref{c1curvewithtangentsonS}, there is a $C^1$ curve $\alpha$ through a further subsequence $\{(f(p_{n_{k}}), \bar{d}f_p(\ell_{n_{k}}))\}$ such that the tangent line of $\alpha$ at each $f(p_{n_{k}})$ is $\bar{d}f_p(\ell_{n_{k}})$.  Thus $f^{-1}(\alpha)$ must be a curve through the points $p_{n_{k}}$ with tangent lines $\ell_{n_{k}}$.  Since this is a subsequence of the transverse sequence $\{(p_n, \ell_n)\}$,  these tangent lines cannot converge.  So $f^{-1}(\alpha)$ cannot be a $C^1$ curve.  But this contradicts the assumption that $f \in \thomeo(S)$.  
\end{proof}


\begin{proof}[Proof of Proposition~\ref{propthm2forward}]
Property $(a)$ comes directly from the definition of $\thomeo(S)$.  The other two desired properties are given by Lemmas~\ref{projtangentspaceishomeo} and \ref{preservetransverse}, respectively.
\end{proof}

\section{Recovering $\bthomeo(\mathbf{S})$}\label{sectionreverse}

The goal of this section is to show that the properties of elements of $\homeo^1(S)$ are sufficient to guarantee that the inverse map must also send $C^1$ curves to $C^1$ curves. More precisely, we prove the following: 

\begin{prop}\label{propthm2reverse}
For a smooth surface $S$, if $f \in \homeo(S)$ has the following three properties:
\begin{enumerate}[noitemsep,topsep=0pt, label=$(\alph*)$]
\item $f$ maps every $C^1$ curve to a $C^1$ curve, 
\item $\bar{d}f_p$ is a homeomorphism for all $p \in S$, and
\item $f$ maps every transverse sequence to a transverse sequence, 
\end{enumerate}
then $f^{-1}$ maps every $C^1$ curve to a $C^1$ curve. 
\end{prop}


We prove this proposition in two steps. We first prove Lemma~\ref{inversetangentline}, showing that every point of the preimage curve has a well-defined tangent line.  We then prove Lemma~\ref{inversetangentcontinuous},  which shows that these tangent lines must vary continuously.

\subsection{Well-defined tangent lines}


We first make a note on what it means for a curve to have a well-defined tangent line at a point.  The tangent line at a point $\gamma(t) = (\gamma_1(t), \gamma_2(t))$ on the parameterized curve $\gamma$ in $\mathbb{R}^2$ is uniquely determined by the slope of the line.  Thus we have a well-defined tangent line if the slopes of the secant lines approaching $\gamma(t)$ have a unique limit.  In other words, there is a well-defined tangent line at $\gamma(t)$ if 
$$\lim_{h \rightarrow 0} \frac{\gamma_2(t+h) - \gamma_2(t)}{\gamma_1(t+h) - \gamma_1(t)}$$ exists.  In our case, we also allow vertical tangent lines if 
$$\liminf_{h \rightarrow 0} \left|\frac{\gamma_2(t+h) - \gamma_2(t)}{\gamma_1(t+h) - \gamma_1(t)} \right| = \infty$$ This notion of a well-defined tangent line can also be translated locally to any surface using a smooth chart around the point $\gamma(t)$. 

With this definition in mind,  the goal of this section is to prove the following: 

\begin{lemma}\label{inversetangentline}
If $f \in \homeo(S)$ preserves tangencies and $\bar{d}f_p$ is a homeomorphism for all $p \in S$, then for any simple closed $C^1$ curve $\gamma$,  the curve $f^{-1}(\gamma)$ has a well-defined tangent line at every point.
\end{lemma}

\noindent Before we prove this lemma,  we note the following observation,  which is a weaker version of Corollary~\ref{c1curvewithtangentsonS}.

\begin{obs}\label{c1connectthedots}
Let $\{x_n\}_{n=1}^{\infty}$ be a sequence of points in $S$ that converge to a point $x$. Then there is a simple closed $C^1$ curve in $S$ that intersects infinitely many of the $x_n$.
\end{obs}

\noindent We can now prove the main result of this section.

\begin{proof}[Proof of Lemma~\ref{inversetangentline}]
Let $\gamma$ be a parameterized $C^1$ curve in $S$ such that $\gamma(0) = f(p)$.  We prove this result by contradiction.  Assume that $f^{-1}(\gamma(0))$ does not have a well-defined tangent line.  Thus after applying a local chart from a neighborhood of $f^{-1}(\gamma(0))$  to $\mathbb{R}^2$,  the limit of the secant slopes does not exist.  Denote the coordinates of $f^{-1}(\gamma(t))$ in the local chart by $(\gamma^{-1}_1(t), \gamma^{-1}_2(t))$.   By considering $\mathbb{R}^2 \cup \{\infty\} \cong S^1$ and using the compactness of  $S^1$,  there must be a slope $m \in \mathbb{R}^2 \cup \{\infty\}$ such that a subsequence of the secant slopes converges to $m$.  More precisely, there exist  $t_{n} \in \left(-\frac 1 n, \frac 1 n \right)$ with
$$\lim_{n \rightarrow \infty} \frac{\gamma^{-1}_2(t_{n}) - \gamma^{-1}_2(0)}{\gamma^{-1}_1(t_{n}) - \gamma^{-1}_1(0)} = m$$
But the limit for all the secant slopes approaching $0$ does not exist, so there is some open neighborhood $N_m$ of $m$ such that there exist $s_n \in \left(-\frac 1 n, \frac 1 n \right)$ with 
$$\qquad \qquad \frac{\gamma^{-1}_2(s_{n}) - \gamma^{-1}_2(0)}{\gamma^{-1}_1(s_{n}) - \gamma^{-1}_1(0)} \notin N_m$$
Thus this sequence of secant slopes is a subset of $S^1 \setminus N_m$, which is also compact.  So there is a subsequence $\{s_{n_k}\}$ such that 
$$\lim_{k\rightarrow \infty}\frac{\gamma^{-1}_2(s_{n_k}) - \gamma^{-1}_2(0)}{\gamma^{-1}_1(s_{n_k}) - \gamma^{-1}_1(0)} = m'$$
For some $m' \neq m$.  By applying a rotation if necessary, we can assume that neither $m$ nor $m'$ are infinity.  

By Observation~\ref{c1connectthedots},  there exist simple closed $C^1$ curves $\alpha$ and $\beta$ such that $\alpha$ intersects infinitely many of the $f^{-1}(\gamma(t_n))$ and $\beta$ intersects infinitely many of the $f^{-1}(\gamma(s_{n_k}))$.  Note that the tangent line of $\alpha$ at $p$ must have a slope of $m$, while the tangent line of $\beta$ at $p$ must have a slope of $m'$.  Since $f(\alpha)$ and $\gamma$ are $C^1$ curves that intersect at infinitely many distinct points that converge to $f(p)$,  then by Corollary~\ref{c1infinitesharetangentline},  $f(\alpha)$ and $\gamma$ must have the same tangent line at $f(p)$. Similarly, $f(\beta)$ and $\gamma$ are $C^1$ curves with the same tangent line at $f(p)$.  Thus $f(\alpha)$ and $f(\beta)$ must have the same tangent line at $f(p)$.  But since $\bar{d}f_p$ is a homeomorphism, this means that $\alpha$ and $\beta$ must have the same tangent line at $p$. This is a contradiction since by construction $m \neq m'$.  
\end{proof}

\subsection{Continuity of tangent lines}

Every $C^1$ curve has tangent lines at every point that vary continuously.  Work of Burgu\'es--Cuf\'i \cite[Theorem 10]{BC} shows that for simple closed curves,  continuously varying tangent lines is sufficient to recover a $C^1$-regular parameterization with nonzero tangent vectors.  

We define a function $M_{\gamma}: \mathbb{R} \rightarrow \mathbb{R} \cup \infty \cong S^1$ that gives the slope the tangent lines of a curve $\gamma$.  Precisely, for a parameterized curve $\gamma(t) = (\gamma_1(t), \gamma_2(t))$ in $\mathbb{R}^2$ this function is defined as
$$M_{\gamma}(t) = \lim_{h \rightarrow 0} \frac{\gamma_2(t+h) - \gamma_2(t)}{\gamma_1(t+h) - \gamma_1(t)}$$
By using a smooth chart, this function is also defined locally for curves on surfaces whenever the tangent line is well-defined.

\begin{lemma}\label{inversetangentcontinuous}
For a surface $S$, let $f \in \homeo(S)$ has all of the following properties:
\begin{enumerate}[noitemsep,topsep=0pt, label=$(\alph*)$]
\item $f$ maps every $C^1$ curve to a $C^1$ curve
\item $\bar{d}f_p$ is a homeomorphism for all $p \in S$
\item $f$ maps every transverse sequence to a transverse sequence.
\end{enumerate}
then for any $C^1$ curve $\gamma$,  $M_{f^{-1}(\gamma)}$ is continuous.
\end{lemma}

\begin{proof}
Let $\gamma(t)$ be a parameterized $C^1$ curve such that $f^{-1}(\gamma(0)) = p$. 
First, note that as a consequence of Lemma~\ref{inversetangentline},  for any sequence $t_n$ in $\mathbb{R}$ that converges to 0,  the sequence $f^{-1}(\gamma(t_n))$ converges along the line with slope $M_{f^{-1}(\gamma)}(0)$ to $p$. By rotating the chart around $p$, we can assume that $M_{f^{-1}(\gamma)}(0)$ is finite.

Now suppose for a contradiction that $M_{f^{-1}(\gamma)}$ is not continuous. Then there is a sequence of $t_n$ in $\mathbb{R}$ that converges to $0$ such that $M_{f^{-1}(\gamma)}(t_n)$ does not converge to $M_{f^{-1}(\gamma)}(0)$.  Moreover, there exists an $\varepsilon >0$ and a subsequence $t_{n_k}$ such that $M_{f^{-1}(\gamma)}(t_{n_k})$ is outside of an $\varepsilon$-neighborhood of $M_{f^{-1}(\gamma)}(0)$. Thus no subsequence of $M_{f^{-1}(\gamma)}(t_{n_k})$ converges to $M_{f^{-1}(\gamma)}(0)$.  Let $p_k$ denote the points $f^{-1}(\gamma)(t_{n_k})$ and let $\ell_k$ denote the element in $\mathbb{P}T_{p_k}S$ that corresponds to $M_{f^{-1}(\gamma)}(t_{n_k})$.  So $\{(p_k,   \ell_k)\}$ is a transverse sequence.  

Note that by construction,  the slope of $\bar{d}f_{p_k}(\ell_k)$  is $M_{\gamma}(t_{n_k})$.  Thus $f$ maps $\{(p_k,   \ell_k)\}$ to a sequence of points and tangent lines from $\gamma$ that converge to $\gamma(0)$.  But $\gamma$ is a $C^1$ curve and thus $M_{\gamma}(t_{n_k})$ converges to $M_{\gamma}(0)$.  Thus $\{f(p_k), \bar{d}f_p(\ell_k)\}$ is not a transverse sequence, contradicting property $(c)$.  Thus $M_{f^{-1}(\gamma)}$ must be continuous.
\end{proof}


\begin{proof}[Proof of Proposition~\ref{propthm2reverse}]
By Lemmas~\ref{inversetangentline} and \ref{inversetangentcontinuous}, $f^{-1}$ maps a $C^1$ curve $\gamma$ to a curve that has continuous tangent line at every point.  The result of Burgu\'es--Cuf\'i \cite[Theorem 10]{BC} implies that $f^{-1}(\gamma)$ is also a $C^1$ curve. 
\end{proof}

\subsection{Finishing the proof}

Combining Propositions~\ref{propthm2forward} and \ref{propthm2reverse}, we now derive our main result.

\begin{proof}[Proof of Main Theorem]
Proposition~\ref{propthm2forward} proves the forward direction of the theorem.  For the reverse direction,  property $(a)$ gives that $f$ maps every $C^1$ curve to a $C^1$ curve.  By Proposition~\ref{propthm2reverse},  $f^{-1}$ also maps every $C^1$ curve to a $C^1$ curve. Thus $f \in \thomeo(S)$. 
\end{proof}

%\bibliographystyle{sn-mathphys-num.bst}
\bibliography{mybib3}



\end{document}
